\section*{Reproducibility appendix}
\label{app:reproducibility}

The software is archived as OTShield Bench \texttt{v0.6.0-alpha} at
\url{https://github.com/ASHTHRA/otshield-bench/releases/tag/v0.6.0-alpha} and
Zenodo DOI \url{https://doi.org/10.5281/zenodo.22899466}. The repository is
\url{https://github.com/ASHTHRA/otshield-bench}. The release declares Python
package version \texttt{0.6.0a0}; the completed study provenance records Python
3.14.4, Docker 29.8.1, Compose v5.5.1, and the Linux/WSL2 kernel used for the
isolated laboratory run.

The repository's software CI runs on Ubuntu and Windows with Python 3.10 and
3.12. POSIX-only v0.6 lifecycle tests are explicitly skipped on native Windows;
portable package, analysis, and accounting tests remain cross-platform. The
manuscript workflow additionally compiles this file and \texttt{main.tex} with
TeX Live in GitHub Actions and uploads the resulting PDF as a workflow artifact.

The study identifier is \texttt{20260922T021855Z}. Its source files are
\texttt{evidence/v06/20260922T021855Z/aggregate.json} and
\texttt{evidence/v06/20260922T021855Z/study.json}; per-run captures and
provenance are under the same evidence root. The release manifest and
\texttt{STUDY\_SHA256SUMS} provide integrity records. The frozen calibration is
\texttt{research/v0.6\_calibration.json}, whose recorded SHA-256 is
\texttt{d8676d4670ba44c0f4b1045e5e07dd2d251ccb4a63ed8089b9b9aada729ea234}.

Software verification is performed with:

\begin{verbatim}
./scripts/verify_release.sh
python paper/generate_figures.py
\end{verbatim}

The first command checks whitespace, runs the complete test suite, and builds
the package. The second reads only committed aggregate data and deterministically
writes the figures in \texttt{paper/figures/}. These commands reproduce the
analysis and manuscript artifacts; they do not launch Docker, OpenPLC, packet
capture, or laboratory traffic.

The figures intentionally read the aggregate rather than duplicating numerical
constants in plotting code. Their condition labels and detector names are the
declared v0.6 schema values; all metric means and interval endpoints come from
\texttt{aggregate.json}. The manuscript's prose is a human-readable rendering
of the same bounded result set and the technical report.

Reproducing the laboratory measurements is a separate activity requiring an
explicitly authorized, isolated Linux/Docker environment. The protocol uses
bounded read-only Modbus/TCP Function Code 3 requests and must never be aimed at
a real operational system. The preserved failed attempts and errata remain in
the repository and are excluded from the completed replacement aggregate.
